\section{Deferred Proofs}\label{app:missing-proofs}

\subsection{Proof of Theorem~\ref{thm:gpd-extension}}\label{app:gpd-extension}
\begin{proof}[Proof of Theorem~\ref{thm:gpd-extension}]
    For a subgraph $H$ of $G$ with $w\br{V_H}\geq B_\eta$, we set $\eta_H=B_\eta/w\br{V_H}$ and $\rho_H=B_\rho/w\br{V_H}$. The algorithm uses $\eta_0=\max\brc{4/5,\eta_H}$ and we observe that $\rho_H/\eta_H=\rho/\eta$.
    Recall that $\OPTdem{H}{\rho_H}$ is the cost of the optimal \textsc{Graph Partitioning with Demands} solution on $H$ with threshold $\rho_H$.
    It is easy to see that the algorithm runs in polynomial time and the returned solution is $\eta$-feasible. Hence, we focus on the approximation ratio.

For a subgraph $H$ of $G$ for which the procedure has been called, let $d\br{H}$ denote the number of recursion levels the algorithm performs on $H$. Every recursive call takes $\eta_0\geq 4/5$ and, by Theorem~\ref{thm:gpd-bicriteria}, produces connected components of demand at most $\eta_0 \cdot w\br{V_H}\leq \frac{4}{5} w\br{V_H}$ each.
Hence, this means that for every `child component' $H'$ of $H\setminus C$ in the recursion tree, we get $\eta_{H'}\geq 5/4\cdot \eta_H$.
Because the last partitioning step happens when $\eta_0=\eta_H\geq 4/5$, the recursion depth is bounded by
\[
    d\br{H}\leq \cl{\log_{5/4}\br{w\br{V_H}/B_\eta}}.
\]
due to the fact that after reaching $\eta_{H}\geq 4/5$ the algorithm performs one last partitioning step, i.e., the last recursive call. Since $\cl{\log_{5/4}\br{w\br{V_H}/B_\eta}}$ is the number of iterations required to reach $\eta_H\geq 1$, which is one more than we need to reach $\eta_H\geq 4/5$, we get that the above quantity is correct.

Now, we prove by induction on $d\br{H}$ that Algorithm~\ref{alg:recursive-partitioning} returns a subset of edges of cost at most $d\br{H}\cdot \kappa\br{\rho,\eta}\cdot \gamma\br{n}\cdot \OPTdem{H}{\rho_H}$.

    \emph{Base cases, $d\br{H}=0$ and $d\br{H}=1$.} If $d\br{H}=0$, then the cost is $0$ since the algorithm terminates without any partitioning. If $d\br{H}=1$, then $B_\eta/w\br{V_H}\geq 4/5$, so $\eta_0=\eta_H=B_\eta/w\br{V_H}$, and after a single partitioning step for each connected component $H'$ of $H\setminus C$, $w\br{V_{H'}}\leq \eta_0 \cdot w\br{V_H}=B_\eta$. Hence the algorithm terminates at the beginning of the next recursive step without adding any supply edges to the cut. By the explicit bound established in the proof of Theorem~\ref{thm:gpd-bicriteria}, the cost of the solution is bounded by $\frac{\gamma\br{n}}{\br{\eta_H+\delta_{\rho_H}-1}^2}\cdot \OPTdem{H}{\rho_H}$. We establish the bound $\frac{1}{\br{\eta_H+\delta_{\rho_H}-1}^2}\leq \kappa\br{\rho,\eta}$ in two cases:
    \begin{enumerate}
        \item If $\rho_H\leq 3/4$, then $\delta_{\rho_H}=1/4$ and $\eta_H+\delta_{\rho_H}-1\geq 4/5+1/4-1=1/20$, using $\eta_H\geq 4/5$. Hence $\frac{1}{\br{\eta_H+\delta_{\rho_H}-1}^2}\leq 400\leq \kappa\br{\rho,\eta}$.
        \item If $\rho_H> 3/4$, then $\delta_{\rho_H}=1-\rho_H$ and, using $\rho_H=\br{\rho/\eta}\cdot\eta_H$ and $\eta_H\geq 4/5$,
        \[
            \eta_H+\delta_{\rho_H}-1=\eta_H-\rho_H=\br{1-\rho/\eta}\cdot\eta_H\geq \br{1-\rho/\eta}\cdot\frac{4}{5}.
        \]
        Hence $\frac{1}{\br{\eta_H+\delta_{\rho_H}-1}^2}\leq \frac{1}{\br{\br{1-\rho/\eta}\cdot 4/5}^2}\leq \kappa\br{\rho,\eta}$.
    \end{enumerate}
    In both cases the cost is at most $\kappa\br{\rho,\eta}\cdot\gamma\br{n}\cdot \OPTdem{H}{\rho_H}=d\br{H}\cdot \kappa\br{\rho,\eta}\cdot\gamma\br{n}\cdot\OPTdem{H}{\rho_H}$.

    \emph{Inductive step.} Suppose the claim holds for every graph $H'$ with $d\br{H'}\leq d-1$, and let $H$ satisfy $d\br{H}=d\geq 2$. Let $C$ be the cut computed when $H$ is the input. It holds $B_\eta/w\br{V_H}<4/5$, so $\eta_0=4/5$, and for every $H'$ in $H\setminus C$, $w\br{H'}\leq \eta_0\cdot w\br{V_H}=\frac{4}{5}\cdot w\br{V_H}$, thus $d\br{H'}\leq d-1$. By Theorem \ref{thm:gpd-bicriteria}, the cost introduced by the current partitioning step is at most $\frac{\gamma\br{n}}{\br{\eta_0+\delta_{\rho_H}-1}^2}\cdot\OPTdem{H}{\rho_H}$, and we again bound $\frac{1}{\br{\eta_0+\delta_{\rho_H}-1}^2}\leq \kappa\br{\rho,\eta}$:
    \begin{enumerate}
        \item If $\rho_H\leq 3/4$, then $\delta_{\rho_H}=1/4$ and $\eta_0+\delta_{\rho_H}-1=4/5+1/4-1=1/20$, so the term is at most $400\leq \kappa\br{\rho,\eta}$.
        \item If $\rho_H> 3/4$, then $\delta_{\rho_H}=1-\rho_H$ and, using $\rho_H=\br{\rho/\eta}\cdot\eta_H<\br{\rho/\eta}\cdot\eta_0$,
        \[
            \eta_0+\delta_{\rho_H}-1=\eta_0-\rho_H> \eta_0-\br{\rho/\eta}\cdot\eta_0=\br{1-\rho/\eta}\cdot\frac{4}{5},
        \]
        so the term is at most $\frac{1}{\br{\br{1-\rho/\eta}\cdot 4/5}^2}\leq \kappa\br{\rho,\eta}$.
    \end{enumerate}
    Thus the top level contributes at most $\kappa\br{\rho,\eta}\cdot\gamma\br{n}\cdot\OPTdem{H}{\rho_H}$. By the inductive hypothesis, for each connected component $H'$ of $H\setminus C$, $c\br{C_{H'}}\leq \br{d-1}\cdot \kappa\br{\rho,\eta}\cdot\gamma\br{n}\cdot\OPTdem{H'}{\rho_{H'}}$. Since the components in $H\setminus C$ are vertex and edge-disjoint and the $C^*\cap E_{H'}$ is a valid solution for \textsc{Graph Partitioning with Demands} on $H'$ in $H\setminus C$ with $\rho_{H'}$, we have that $\OPTdem{H'}{\rho_{H'}}\leq c\br{C^*\cap E_{H'}}$. Moreover, by summing up over all $H'$ we get
    $\sum_{H'\in H\setminus C}\OPTdem{H'}{\rho_{H'}}\leq \OPTdem{H}{\rho_H}$, and therefore
    \begin{align*}
        c\br{C}&=c\br{C_H}+\sum_{H'\in H\setminus C}c\br{C_{H'}}
        \\&\leq \kappa\br{\rho,\eta}\cdot\gamma\br{n}\cdot\OPTdem{H}{\rho_H}+\br{d-1}\cdot\kappa\br{\rho,\eta}\cdot \gamma\br{n}\sum_{H'\in H\setminus C}\OPTdem{H'}{\rho_{H'}}
        \\&\leq 
        d\cdot \kappa\br{\rho,\eta}\cdot\gamma\br{n}\cdot\OPTdem{H}{\rho_H}.
    \end{align*}
    Therefore, the induction is complete. Applying it to $H=G$ yields the statement of the theorem.
\end{proof}

\subsection{Proof of Lemma~\ref{lem:tree-separator}}\label{app:tree-separator}
\begin{proof}[Proof of Lemma~\ref{lem:tree-separator}]
    Root $H$ arbitrarily and initialize the current root-containing subtree $R$ to $H$. While $w\br{V_R}>w\br{V_H}/t$, choose the deepest vertex $x$ such that the subtree $R_x$ rooted at $x$ in the current $R$ satisfies $w\br{V_{R_x}}>w\br{V_H}/t$.
    Add $x$ to $\mathcal L_H$ and remove $R_x$ from $R$. The peeled subtrees are pairwise disjoint, so the sum of their internal demand weights is at most $w\br{V_H}$. Since every peeled subtree has demand greater than $w\br{V_H}/t$, there are at most $t-1$ of them. By the choice of a deepest such vertex, every child subtree $T'$ of $x$ in $R_x$ satisfies $w\br{V_{T'}}\leq w\br{V_H}/t$. Therefore, after removing all separators, every connected component $K$ in $H\setminus\mathcal L_H$ is one of these child subtrees at the iteration when its parent separator was selected, or the final remainder $R$. The final remainder satisfies $w\br{V_R}\leq w\br{V_H}/t$ by the stopping condition. In either case, $w\br{V_K}\leq w\br{V_H}/t$ as desired.
\end{proof}

\subsection{Cost bound in Theorem~\ref{thm:tree-centroid-max-coverage}}\label{app:tree-centroid-max-coverage-cost}
\begin{proof}[Proof of the cost bound in Theorem~\ref{thm:tree-centroid-max-coverage}]
Let $\OPT_{B_\rho}(H)$ denote the minimum cost of a cut that leaves every component of $H$ with internal demand at most $B_\rho$. We prove by induction on the maximum number $d(H)$ of nonterminal recursive calls on a path starting at $H$ that the total cost $C_H$ returned below $H$ satisfies
\[
    c(C_H)\leq d(H)\cdot\OPT_{B_\rho}(H).
\]
For a terminal component, $d(H)=0$ and $C_H=\emptyset$, so the claim holds. Otherwise, let $F_H$ be the cut selected at $H$, and let $H_1,\ldots,H_m$ be the components of $H\setminus F_H$. Then
\[
    d(H)=1+\max_{1\leq i\leq m}d(H_i).
\]

Let $F_H^*$ be an optimal $B_\rho$-feasible cut of $H$, and set $Q^*=F_H^*\cap E_{S_H}$. The algorithm considers this guess. For each $K\in\mathcal K_{Q^*}^+$, the component of $K\setminus F_H^*$ containing $S_K$ has internal demand at most $B_\rho$. Hence the restriction of $F_H^*$ to $K$ separates at least $q_K=w(V_K)-B_\rho$ demand from $S_K$. It follows that
\[
    \OPT_{q_K}(K)\leq c(F_H^*\cap E_K).
\]
Since $\textsc{CoreQuotaCover}$ returns a set $A_K$ with $f_K(A_K)\geq(1-1/e)q_K$ and $c_K(A_K)\leq\OPT_{q_K}(K)$, the candidate for $Q^*$ has cost at most
\[
    c(Q^*)+\sum_{K\in\mathcal K_{Q^*}^+}\OPT_{q_K}(K)
    \leq c(F_H^*)=\OPT_{B_\rho}(H).
\]
Since the algorithm chooses the least-cost candidate, $c(F_H)\leq\OPT_{B_\rho}(H)$.

For each child $H_i$, restricting $F_H^*$ to $E_{H_i}$ is feasible for $\OPT_{B_\rho}(H_i)$. The child edge sets are pairwise disjoint, so
\[
    \sum_{i=1}^m\OPT_{B_\rho}(H_i)
    \leq\sum_{i=1}^m c(F_H^*\cap E_{H_i})
    \leq\OPT_{B_\rho}(H).
\]
Assuming the induction claim for each child, we obtain
\begin{align*}
    c(C_H)
    &=c(F_H)+\sum_{i=1}^m c(C_{H_i})\\
    &\leq\OPT_{B_\rho}(H)+(d(H)-1)\cdot \sum_{i=1}^m\OPT_{B_\rho}(H_i)\\
    &\leq d(H)\cdot\OPT_{B_\rho}(H).
\end{align*}
By the excess-demand bound in the proof of Theorem~\ref{thm:tree-centroid-max-coverage}, all recursive calls terminate by depth
\[
    L=\left\lceil\ln\left(\frac{1-\rho}{\eta-\rho}\right)\right\rceil.
\]
Thus $d(T)\leq L$, and the claimed cost bound follows.
\end{proof}

\subsection{Proof of Theorem~\ref{thm:tree-centroid-full-quota}}\label{app:tree-centroid-full-quota}
\begin{proof}[Proof of Theorem~\ref{thm:tree-centroid-full-quota}]
    By Lemma~\ref{lem:tree-separator}, $\spr{\mathcal{L}_T}\leq t-1$, and for every $H\in T\setminus\mathcal L_T$, $w\br{H}\leq W/t\leq B_\rho$. By construction, all components $H\in T\setminus C$ have demand bounded by $w\br{V_H}\leq B_\rho$.

    Let $F^*$ be an optimal $B_\rho$-feasible cut and set $Q^*=F^*\cap E_{S_T}$. For each $K\in\mathcal K_{Q^*}^+$, fix an optimal $A_K^*$ for $\OPT_{q_K}(K)$. To obtain a valid multicut $M$ for $T_K$, for each edge $e=uv\in A_K^*$ include both representative edges $u^{\mathrm L}v^{\mathrm L}$ and $u^{\mathrm R}v^{\mathrm R}$ in $M$. This doubles the cost of each edge, so $c(M)\leq2\OPT_{q_K}(K)$. Mestre's $(2+\delta)$-approximation returns a multicut $\widehat A_K$ with
    \[
        c(\widehat A_K)\leq(2+\delta)\cdot c(M)\leq(4+\epsilon)\cdot\OPT_{q_K}(K),
    \]
    where $\delta=\epsilon/2$. The routine returns $A_K$ by replacing every selected copied edge in $\widehat A_K$ by its corresponding original supply edge and merging duplicates. This preserves the separated demand and cannot increase cost, so $c(A_K)\leq c(\widehat A_K)\leq(4+\epsilon)\cdot\OPT_{q_K}(K)$.

    The restrictions of $F^*$ to the disjoint components in $\mathcal K_{Q^*}^+$ are feasible for their quota problems, and their edges together with $Q^*$ are a subset of $F^*$. Therefore,
    \[
        c(Q^*)+\sum_{K\in\mathcal K_{Q^*}^+}\OPT_{q_K}(K)\leq\OPTdem{T}{\rho}.
    \]
    Since the algorithm considers $Q^*$ and chooses a minimum-cost candidate, its output satisfies
    \[
        c(C)\leq c(Q^*)+\sum_{K\in\mathcal K_{Q^*}^+}c(A_K)
        \leq(4+\epsilon)\cdot\OPTdem{T}{\rho}.
    \]

    Since the number of sets $Q$ we enumerate is bounded by $\sum_{i=0}^{t-1} \binom{|E_{S_T}|}{i}=n^{O(1/\rho)}$, the algorithm runs in polynomial time for any fixed choice of $\rho$.
\end{proof}
